% !TeX root = ../../Book3.tex
\chapter{Riemann 曲面与代数曲线基础}
\label{b3:app:A}
\begin{chapterabstract}
复平面的坐标可以沿解析延拓改变，也可以由多项式方程隐式给出。本附录把这些局部坐标组织成 Riemann 曲面，说明亚纯函数、微分形式、分歧覆盖和除子怎样相互联系。Riemann--Hurwitz 公式用局部重数计算拓扑亏格；Riemann--Roch 公式则计算给定极点与零点条件下的全局函数空间。从坐标、主部与拼接障碍出发，可以看清维数公式如何给出亚纯函数的存在性和插值，进而得到射影嵌入及亏格一曲面的周期表示。
\end{chapterabstract}
\begin{learninggoals}
\begin{itemize}
\item\label{b3:item:A-goal-1}区分坐标函数、全局函数与全局微分形式，计算换坐标和拉回后的阶数及留数。
\item\label{b3:item:A-goal-2}从局部标准形数出覆盖层数，用分歧数据计算亏格并检查无穷远点。
\item\label{b3:item:A-goal-3}用除子表达极点与零点条件，辨认局部主部的自由度和全局拼接的限制。
\item\label{b3:item:A-goal-4}识别 Riemann--Roch 的深层输入，正确使用其插值与嵌入推论，并区别平面曲线与光滑模型。
\end{itemize}
\end{learninggoals}

在平面区域内，函数与导数都用同一个变量表达。把一族解析分支接起来以后，通常找不到覆盖所有点的单一变量。例如平方根的两层在根处相接，环面上的平面坐标又会相差一个周期。此时需要先确定哪些表达只属于一个坐标图，哪些对象能够在换坐标后保持同一意义。微分形式、除子和层分别记录转换规律、局部阶数以及相容拼接的条件。

\ProseParagraphBreak
局部零点的阶数在换坐标后保持不变，因此可以用它记录映射的分歧，也可以用它给函数施加全局的零极点条件。Riemann--Hurwitz 公式把分歧转成亏格，Riemann--Roch 公式则把零极点条件转成函数空间的维数；后一个公式只给证明概要，所需深层结果与参考文献在相应定理旁说明。

\section{复结构与局部映射}
\label{b3:sec:A01}

\subsection{坐标相容与全局对象}
\label{b3:subsec:A0101}

局部坐标的任务是把一小片空间变成平面开集；两张图的重叠部分必须对“全纯”给出相同判断。因此，决定复结构的是坐标之间的转换，而非把某一个坐标选为全局变量。

\Needspace{6\baselineskip}
\begin{definition}[Riemann 曲面]\label{b3:def:A-riemann-surface}
设 \(X\) 是连通、Hausdorff、第二可数的拓扑空间。若存在开覆盖 \((U_j)\) 及同胚 \(z_j:U_j\rightarrow V_j\)，其中 \(V_j\subset\mathbb C\) 开，而且在每个非空重叠处，坐标转换 \(z_k\circ z_j^{-1}\) 都双全纯，则称带此复结构的 \(X\) 为\term[riemann-qumian]{Riemann 曲面}[Riemann surface]。这样的坐标族称为复图册；两套图册的所有交叉转换均双全纯时，称它们相容，并视为确定同一复结构。
\end{definition}

Hausdorff 条件使不同点可由不交邻域分离，第二可数条件指拓扑有一个可数基。后者便于用可数坐标图进行局部构造；二者不能由“每一点看起来像平面”自动代替。平面区域的恒等坐标给出最简单的例子。Riemann 球面 \(\widehat{\mathbb C\kern0pt}\) 用有限坐标 \(z\) 和无穷远坐标 \(t=1/z\) 覆盖；重叠处的转换导数 \(-1/z^2\) 不为0。

\ProseParagraphBreak
再取二维实格点群 \(\Lambda=\mathbb Z\lambda_1+\mathbb Z\lambda_2\)，其中 \(\lambda_1,\lambda_2\) 实线性无关。在 \(\mathbb C/\Lambda\) 中，一个充分小的平面圆盘不与它的非平凡格点平移相交，商映射在其上单射，反函数就是商空间的坐标。两种提升的差为格点，转换形如 \(z\mapsto z+\lambda\)。平面函数 \(z\) 因此不能下降为环面上的复值函数，而它的微分 \(\mathrm{d}z\) 可以下降：平移不改变微分。

\ProseParagraphBreak
复图册还给一个相容的实定向。若转换为 \(w=\varphi(z)\)，Cauchy--Riemann 方程给
\[
 \det D_{\mathbb R}\varphi=|\varphi'(z)|^2>0.
\]
所以在不同坐标中判定的正向小圆和正向面积相符。后面的留数积分与 Euler 特征都使用这一自然定向。

\subsection{局部重数与标准坐标}
\label{b3:subsec:A0102}

连续映射 \(F:X\rightarrow Y\) 称为全纯，是说在源、目标坐标中，每个表达 \(w\circ F\circ z^{-1}\) 都全纯。只需在覆盖各点的一套图中检查，因为换图仅在左右复合双全纯映射。平面恒等定理也适用于曲面：局部相等沿相交坐标圆盘传播，连通性使相等延伸到整个曲面。特别地，非常值全纯映射不能在某个开集上为常值。

\Needspace{6\baselineskip}
\begin{proposition}[全纯映射的局部标准形]\label{b3:prop:A-local-normal-form}
设 \(X,Y\) 为 Riemann 曲面，\(F:X\rightarrow Y\) 为非常值全纯映射，\(P\in X\)、\(Q=F(P)\)。选以 \(Q\) 为原点的目标坐标 \(w\)，则存在以 \(P\) 为原点的源坐标 \(t\) 及唯一整数 \(e\ge1\)，使
\[
 w(F(P'))=t(P')^e
\]
在 \(P\) 的一个邻域成立。整数 \(e\) 与两边的坐标选择无关，称为 \(F\) 在 \(P\) 的\term[jubu-chongshu]{局部重数}[local multiplicity]，记为 \(e_P(F)\)。当且仅当 \(e_P(F)=1\) 时，\(F\) 在 \(P\) 附近双全纯。
\end{proposition}
\begin{proof}
先取任意源坐标 \(z(P)=0\)。目标坐标表达不是恒零函数，故其 Taylor 展开有首个非零项，可写为 \(z^e u(z)\)，其中 \(u(0)\ne0\)。缩小圆盘使 \(u\) 无零点，并在圆盘中选全纯的 \(e\) 次根 \(v\)。令 \(t=zv(z)\)，则 \(t'(0)=v(0)\ne0\)，由\cref{b3:thm:01-local-inverse}，\(t\) 是新的局部坐标，而原表达等于 \(t^e\)。

若换源坐标，新旧参数之间的转换具有一阶零点，代入不会改变零点阶数。若换以 \(Q\) 为原点的目标坐标，转换形如 \(w\mapsto w a(w)\)，\(a(0)\ne0\)，也不改变阶数。因此 \(e\) 不依赖坐标。\(e=1\) 时标准形本身可逆；\(e>1\) 时任意充分小的非零目标值都有 \(e\) 个不同局部原像，因而没有包含 \(P\) 的单射邻域。
\end{proof}

若 \(e_P(F)>1\)，称 \(P\) 为\term[fenqi-dian]{分歧点}[ramification point]；其像称为\term[fenqi-zhi]{分歧值}[branch value]。前者在源曲面，后者在目标曲面。局部式 \(t^e\) 还说明非常值全纯映射是开映射，以及一个分歧点附近的 \(e\) 层怎样在中心合并。这里数的是该小坐标邻域内的原像，尚未计算整个曲面上的总原像数。

\ProseParagraphBreak
亚纯函数的极点可用相同语言处理。若 \(f\) 在 \(P\) 有 \(m\) 阶极点，取球面上以无穷远为中心的目标坐标 \(w=1/z\)，则 \(w\circ f=1/f\) 有 \(m\) 阶零点。所以亚纯函数恰是到球面的全纯映射，极点处的局部重数就是极点阶数；极点并不是这个球面值映射的不连续点。

\subsection{紧性对全局映射的限制}
\label{b3:subsec:A0103}

\Needspace{6\baselineskip}
\begin{proposition}[紧曲面的函数]\label{b3:prop:A-compact-functions}
紧 Riemann 曲面上的全纯复值函数都是常数。若 \(X,Y\) 是 Riemann 曲面，\(X\) 紧且 \(F:X\rightarrow Y\) 为非常值全纯映射，则 \(F\) 满射，\(Y\) 也紧，并且每条纤维 \(F^{-1}(Q)\) 都是非空有限集。
\end{proposition}
\begin{proof}
全纯复值函数的模在紧集上达到最大值。最大点的坐标圆盘内可用最大模原理，函数在此为常数，再由曲面上的恒等定理得全局常数。

非常值 \(F\) 由局部标准形是开映射。像 \(F(X)\) 又是紧集，在 Hausdorff 空间 \(Y\) 中闭；\(Y\) 连通，故非空的开闭子集 \(F(X)\) 等于 \(Y\)，也证明 \(Y\) 紧。每条纤维是 \(X\) 的闭子集，因而紧；局部标准形使它离散。用每个点只含该纤维一点的邻域覆盖纤维，紧性给有限子覆盖，故纤维有限。
\end{proof}

这解释了为什么紧曲面的函数理论必须允许极点。球面上的 \(z\) 和环面上的椭圆函数都是非常值亚纯函数；一般紧曲面也有这种函数，但其存在性将在 Riemann--Roch 之后证明。不能仅因每一点有局部坐标，就断言这些坐标能够拼成一个非常值全局函数。

\section{微分形式、留数与除子}
\label{b3:sec:A02}

\subsection{微分形式的转换与拉回}
\label{b3:subsec:A0201}

在坐标 \(z\) 中，1形式写为 \(\omega=a(z)\,\mathrm{d}z\)。它要能沿曲线积分，换成坐标 \(w\) 后就必须遵守链式法则：
\[
 \omega=a(z(w))\cdot\frac{\mathrm{d}z}{\mathrm{d}w}\,\mathrm{d}w.
\]
各图中的系数按这一规律相容，就得到曲面上的1形式。系数全纯或亚纯时，分别称为全纯或亚纯1形式。系数不变并非拼接的要求；真正需要不变的是系数乘微分的整个表达。

\ProseParagraphBreak
例如球面上的 \(\mathrm{d}z\) 在无穷远坐标 \(t=1/z\) 中是 \(-t^{-2}\,\mathrm{d}t\)，有二阶极点。环面上的 \(\mathrm{d}z\) 则因坐标转换为平移而成为处处无零的全纯形式。对于非零亚纯1形式，定义 \(\operatorname{ord}_P\omega\) 为局部系数的零点阶数或负的极点阶数。坐标转换导数处处非零，所以阶数与坐标无关。

\ProseParagraphBreak
\(\operatorname{Res}_P\omega\) 是以 \(P\) 为原点的局部 Laurent 展开中 \(z^{-1}\,\mathrm{d}z\) 的系数，也等于绕 \(P\) 的正向小圈积分除以 \(2\pi\mathrm{i}\)。换坐标给同一条曲线上的积分；变形到另一个正向小圈不改变该积分，因而留数也与坐标无关。若 \(F:X\rightarrow Y\) 全纯，则 \(F^*\omega\) 在坐标中由代入和链式法则定义，称为形式的拉回。

\Needspace{6\baselineskip}
\begin{proposition}[拉回的阶数与留数]\label{b3:prop:A-pullback-differential}
设 \(F:X\rightarrow Y\) 为 Riemann 曲面间的非常值全纯映射，\(P\in X\)、\(Q=F(P)\)、\(e=e_P(F)\)。若 \(\omega\) 是 \(Q\) 附近非零的亚纯1形式，\(m=\operatorname{ord}_Q\omega\)，则
\[
 \operatorname{ord}_P(F^*\omega)=e(m+1)-1,
 \quad
 \operatorname{Res}_P(F^*\omega)=e\operatorname{Res}_Q\omega.
\]
\end{proposition}
\begin{proof}
用局部标准形 \(w=t^e\)，写 \(\omega=\sum_{k=m}^{\infty}a_k w^k\,\mathrm{d}w\)，其中 \(a_m\ne0\)。则
\[
 F^*\omega=\sum_{k=m}^{\infty}e a_k t^{e(k+1)-1}\,\mathrm{d}t.
\]
首项指数为 \(e(m+1)-1\)。指数等于 \(-1\) 当且仅当 \(k=-1\)，相应系数为 \(e a_{-1}\)，这同时给出两式。
\end{proof}

当目标形式在 \(Q\) 全纯且不为0时，拉回在 \(P\) 的零点阶数为 \(e-1\)；分歧本身就产生零点。当目标形式有简单极点时，拉回仍有简单极点，留数却乘以 \(e\)。因此拉回极点阶数不能只照搬函数阶数的乘法规则，微分中的导数还贡献 \(e-1\)。

\subsection{局部原函数、周期与全局留数和}
\label{b3:subsec:A0202}

在穿孔坐标圆盘内，逐项积分 Laurent 展开时，\(a_k z^k\,\mathrm{d}z\) 对 \(k\ne-1\) 给 \(a_k z^{k+1}/(k+1)\)，只有 \(a_{-1}\,\mathrm{d}z/z\) 需要多值对数。因此亚纯1形式在此有单值原函数，当且仅当留数为0；这时原函数也亚纯。这个局部结论没有检查绕曲面其他闭路的积分。

\ProseParagraphBreak
更确切地，设紧曲面 \(X\) 上的亚纯形式 \(\omega\) 的极点集为 \(S\)。它有全局亚纯原函数，当且仅当
\[
 \int_\gamma\omega=0
\]
对 \(X\setminus S\) 中所有分段光滑闭路 \(\gamma\) 成立。必要性是闭路上微分积分为0。反之，\(X\setminus S\) 仍路径连通：任意路径经过有限个被删点时，可在小坐标盘中绕开。固定基点后以路径积分定义 \(F\)，闭路积分条件保证路径无关，局部积分给 \(\mathrm{d}F=\omega\)。每个极点附近的小圈周期为0，所以刚才的局部原函数可将 \(F\) 亚纯延拓到该点。闭路积分称为形式的周期；环面上的 \(\mathrm{d}z\) 没有极点，却沿两条生成闭路有非零周期，故留数条件不能代替全部周期条件。

\Needspace{6\baselineskip}
\begin{theorem}[紧曲面的留数和]\label{b3:thm:A-global-residue}
设 \(X\) 为紧 Riemann 曲面，\(\omega\) 为 \(X\) 上亚纯1形式。则其极点有限，并有
\[
 \sum_{P\in X}\operatorname{Res}_P\omega=0.
\]
\end{theorem}
\begin{proof}
每点都有系数亚纯的坐标邻域，极点集局部离散；有限个这样的邻域覆盖紧曲面，故极点有限。在各极点周围挖去互不相交的小坐标圆盘。使用紧光滑曲面的有限相容三角剖分，并充分细分，使余下区域可分成各自位于一个坐标图内的有限片；此处采用三角剖分的标准存在性作为拓扑输入。

每片内形式全纯，由平面 Cauchy 定理或 Stokes 公式，沿该片边界的积分为0。相邻片的共同边反向出现，积分抵消；形式的转换规律保证两边算的是同一个积分。剩下的是挖孔边界，它们相对于余下曲面都顺时针定向。因此各正向小圈积分之和为0，再除以 \(2\pi\mathrm{i}\) 得留数和公式。
\end{proof}

例如在环面上不可能存在只有一个简单极点的亚纯函数 \(f\)。若存在，用处处无零的形式 \(\mathrm{d}z\) 相乘，\(f\,\mathrm{d}z\) 将只有一个简单极点且留数非零，违背上式。这是全局拓扑限制局部主部的第一个例子。

\subsection{除子与线性等价}
\label{b3:subsec:A0203}

记录零点和极点的位置及阶数，不需要保留函数的每个 Taylor 系数。把点作为形式符号，就可以将这些离散数据相加，并把阶数不等式写成一个整体条件。

\Needspace{6\baselineskip}
\begin{definition}[除子]\label{b3:def:A-divisor}
设 \(X\) 为紧 Riemann 曲面。\(X\) 上的\term[chuzi]{除子}[divisor]是有限形式和 \(D=\sum_P n_PP\)，其中 \(n_P\in\mathbb Z\)。记 \(D(P)=n_P\)，其\term[chuzi-cishu]{次数}[degree of a divisor]为 \(\deg D=\sum_P n_P\)。所有系数非负时称 \(D\) 为\term[youxiao-chuzi]{有效除子}[effective divisor]，记为 \(D\ge0\)。非零亚纯函数的\term[zhuchuzi]{主除子}[principal divisor]为
\[
 (f)=\sum_P\operatorname{ord}_P(f)P.
\]
若 \(D'-D=(h)\)，其中 \(h\) 是非零亚纯函数，则称 \(D,D'\) \term[xianxing-dengjia]{线性等价}[linearly equivalent]，记为 \(D\sim D'\)。非零亚纯1形式同样定义除子 \((\omega)=\sum_P\operatorname{ord}_P(\omega)P\)。
\end{definition}

\symidx{\(\deg D\)：除子的次数}
这里的和有限，因为非零亚纯函数或形式的零极点局部离散，而曲面紧。若 \(f=z^m u(z)\)、\(u(0)\ne0\)，则
\[
 \frac{\mathrm{d}f}{f}=m\frac{\mathrm{d}z}{z}+\frac{\mathrm{d}u}{u}.
\]
右侧留数为 \(m\)。对全局形式 \(\mathrm{d}f/f\) 用留数和定理，得到 \(\deg(f)=0\)：非零亚纯函数的零点总阶数与极点总阶数相等。于是线性等价保持次数。若两个函数的除子相同，它们的商无零点、无极点，紧性使此商为非零常数。

\ProseParagraphBreak
次数相同仍不保证线性等价。环面上不同点 \(P,Q\) 的除子 \(P-Q\) 次数为0，却不是主除子；否则实现它的函数在 \(Q\) 只有一个简单极点，违反上一小节的留数论证。相比之下，球面上有限点 \(a,b\) 满足
\[
 \left(\frac{z-a}{z-b}\right)=[a]-[b].
\]
这里方括号只是强调该符号表示一个点；无穷远的贡献已经在分子、分母中抵消。

\subsection{允许极点的有限维空间}
\label{b3:subsec:A0204}

给定紧曲面上的除子 \(D\)，定义
\[
 L(D)\coloneqq\{0\}\cup
 \{\,f\text{ 为非零亚纯函数}\mid(f)+D\ge0\,\},
 \quad \ell(D)\coloneqq\dim_{\mathbb C}L(D).
\]
\symidx{\(L(D)\)：允许除子指定极点的亚纯函数空间}\symidx{\(\ell(D)\)：空间 \(L(D)\) 的维数}
若 \(D(P)=m\)，局部条件就是 \(t^m f\) 全纯：\(m>0\) 时允许至多 \(m\) 阶极点，\(m<0\) 时要求至少 \(-m\) 阶零点。这种表达也立即证明 \(L(D)\) 是复向量空间，因为相加不会降低两个局部函数阶数的较小值。

\ProseParagraphBreak
在球面上，任意亚纯函数都是有理函数。确实，减去有限极点的有限 Laurent 主部，再减去无穷远主部对应的多项式，余下的函数在整个球面全纯，故为常数。由此，若 \(0\le m\le n\)，则
\[
 L(n[\infty]-m[0])
 =\operatorname{span}_{\mathbb C}\{z^m,z^{m+1},\ldots,z^n\}.
\]
这里要求在0消失 \(m\) 阶，同时在无穷远至多有 \(n\) 阶极点。若 \(m>n\)，空间只含0。

\Needspace{6\baselineskip}
\begin{proposition}[允许极点空间的有限维性]\label{b3:prop:A-divisor-finite}
设 \(X\) 为紧 Riemann 曲面，\(D=D_+-D_-\) 为其除子，\(D_+,D_-\) 有效且支集不交。定义 \(L(D)=\{0\}\cup\{f:(f)+D\ge0\}\)，其中非零 \(f\) 为 \(X\) 上亚纯函数，\(\ell(D)=\dim_{\mathbb C}L(D)\)。则
\[
 \ell(D)\le\ell(D_+)\le1+\deg D_+.
\]
\end{proposition}
\begin{proof}
第一步来自 \(L(D)\subset L(D_+)\)。写 \(D_+=\sum_{j=1}^s m_jP_j\)，在各 \(P_j\) 选坐标，在这些点以外另选一点 \(Q\)。把 \(f\in L(D_+)\) 映到各 \(P_j\) 的 \(m_j\) 个负次 Laurent 系数及数值 \(f(Q)\)，得到至 \(\mathbb C^{1+\sum m_j}\) 的线性映射。若这些数据全为0，所有可能的极点都可去，\(f\) 是紧曲面上的全纯函数，因而为常数；\(f(Q)=0\) 又使此常数为0。映射单射，维数上界随之成立。
\end{proof}

几个直接事实将在最后一节反复使用。若 \(\deg D<0\)，则 \(L(D)=\{0\}\)，否则有效除子 \((f)+D\) 的次数将为负。若 \(\deg D=0\) 且有非零 \(f\in L(D)\)，则有效除子 \((f)+D\) 次数为0，必为零除子，故 \(D=-(f)\) 为主除子；所有这样的函数彼此相差常数因子。因此次数0时，\(\ell(D)=1\) 当且仅当 \(D\) 为主除子，其余情形 \(\ell(D)=0\)。

\ProseParagraphBreak
若 \(D'=D+(h)\)，映射 \(f\mapsto f/h\) 给 \(L(D)\) 到 \(L(D')\) 的线性同构，因为 \((f/h)+D'=(f)+D\)。最后，对任意点 \(P\) 和整数 \(m=D(P)\)，局部系数映射
\[
 L(D+P)\longrightarrow\mathbb C,
 \quad f\longmapsto(t^{m+1}f)(P)
\]
的核恰为 \(L(D)\)，包括 \(m<0\) 的情形。因此
\[
 0\le\ell(D+P)-\ell(D)\le1.
\]
多允许一个局部系数至多增加一维，但这个系数是否能够任意取值，要由全局条件决定。Riemann--Roch 将给出这些限制的总量。

\section{分歧覆盖与拓扑计数}
\label{b3:sec:A03}

\subsection{纤维重数与映射次数}
\label{b3:subsec:A0301}

局部标准形只描述一小片源曲面。要把附近的根数拼成一个全局次数，还需保证目标点移动时没有其他原像从未控制的区域进入。这正是紧性在有限覆盖论证中的作用。

\Needspace{6\baselineskip}
\begin{proposition}[纤维的重数和]\label{b3:prop:A-degree-fibers}
设 \(F:X\rightarrow Y\) 为紧 Riemann 曲面之间的非常值全纯映射。它的分歧点和分歧值均有限；删去全部分歧值 \(B\) 及其原像后，\(F:X\setminus F^{-1}(B)\rightarrow Y\setminus B\) 是有限层覆盖。存在整数 \(d\ge1\)，使每个 \(Q\in Y\) 都满足
\[
 \sum_{P\in F^{-1}(Q)}e_P(F)=d.
\]
这个 \(d\) 称为 \(F\) 的\term[yingshe-cishu]{次数}[degree of a map]。若 \(d=1\)，则 \(F\) 双全纯。
\end{proposition}
\begin{proof}
在坐标中分歧点是映射导数的零点。导数不恒为0，故这些点局部离散，紧性使其总数有限。

固定 \(Q\)，由\cref{b3:prop:A-compact-functions}，其原像为有限集 \(P_1,\ldots,P_s\)。在这些点周围取互不相交的标准形坐标盘 \(U_j\)，并使 \(w\circ F=t_j^{e_j}\)。紧集 \(K=X\setminus\bigcup_j U_j\) 的像不含 \(Q\)，所以可把目标坐标盘缩小至避开 \(F(K)\)。再缩小半径，使它在每张标准图内的全部 \(e_j\) 个根都留在 \(U_j\)。这样，小目标盘的所有原像均在这些图中，没有额外原像层。

若 \(Q\) 非分歧值，则所有 \(e_j=1\)，各图给目标小盘的一张双全纯原像层，证明覆盖性质。曲面删去有限点仍连通，因此该覆盖的层数是一个固定整数 \(d\)。对于分歧值 \(Q\)，取其小盘中的正则值；第 \(j\) 个标准图贡献 \(e_j\) 个原像，故 \(\sum_j e_j=d\)。若 \(d=1\)，每个纤维只有一点，且每处重数都为1；局部全纯逆拼成全局全纯逆。
\end{proof}

局部重数在复合下相乘：将 \(u=t^e\)、\(w=u^r\) 复合得到 \(w=t^{er}\)。若 \(F:X\rightarrow Y\)、\(G:Y\rightarrow Z\) 都是紧曲面间的非常值全纯映射，对 \(Z\) 的一点逐层求重数和，即得 \(\deg(G\circ F)=\deg G\cdot\deg F\)。特别地，一个亚纯函数作为到球面的映射，其次数等于极点总阶数，也等于任意有限值的原像重数和。

\subsection{亏格与 Euler 特征}
\label{b3:subsec:A0302}

本节沿用\subsecref{b3:subsec:120303}的拓扑亏格：闭定向曲面的柄数。所用拓扑输入是曲面分类、与光滑结构相容的有限三角剖分，以及可使指定有限点集成为顶点的细分。它们的存在性不在本附录证明。有限胞腔计数的交错和
\[
 \chi(X)=V-E+T
\]
与分解选择无关，称为 Euler 特征；这里 \(V,E,T\) 分别为点、开边和开面的数目。此种计数不变性可参见 Hatcher 的《Algebraic Topology》\cite[pp. 146--147, Theorem 2.44]{Hatcher2002Algebraic}。

\ProseParagraphBreak
对于 \(g\ge1\)，标准亏格 \(g\) 曲面由一个 \(4g\) 边形按边界词
\[
 a_1b_1a_1^{-1}b_1^{-1}\cdots a_gb_ga_g^{-1}b_g^{-1}
\]
粘合而成。同名字的边配对，逆号表示反向识别；所有顶点最终合成一点，\(4g\) 条边留下 \(2g\) 条开边，内部留下一个开面。因此 \(\chi(X)=1-2g+1=2-2g\)。球面可由闭圆盘把整个边界压为一点，得到一个点和一个开面，仍有 \(\chi=2\)。这些模型见同书\cite[pp. 5--6, 141]{Hatcher2002Algebraic}。随后换用细剖分时不必重新识别柄数，因为交错和保持不变。

\subsection{Riemann--Hurwitz 与分歧除子}
\label{b3:subsec:A0303}

无分歧的 \(d\) 层覆盖把每个小面、每条边和每个顶点都复制 \(d\) 份。分歧只在有限点处把若干原像合并；\term[riemann-hurwitz-gongshi]{Riemann--Hurwitz 公式}[Riemann--Hurwitz formula]计算的正是这种顶点缺额。

\Needspace{6\baselineskip}
\begin{theorem}[Riemann--Hurwitz 公式]\label{b3:thm:A-riemann-hurwitz}
设 \(X,Y\) 是亏格分别为 \(g_X,g_Y\) 的紧 Riemann 曲面，\(F:X\rightarrow Y\) 为次数 \(d\ge1\) 的非常值全纯映射，\(e_P(F)\) 为其在 \(P\) 的局部重数。则
\[
 2g_X-2=d(2g_Y-2)+\sum_{P\in X}\bigl(e_P(F)-1\bigr).
\]
\end{theorem}
\begin{proof}
在 \(Y\) 上取有限相容三角剖分，使所有分歧值都成为顶点。每条开边与每个开面避开分歧值，并且单连通，故各有 \(d\) 个提升。若目标剖分的顶点、边、面数为 \(V,E,T\)，提升的边、面数为 \(dE,dT\)。在顶点 \(Q\) 上方，由纤维重数和公式，点数为
\[
 \#F^{-1}(Q)=d-\sum_{F(P)=Q}(e_P(F)-1).
\]
所以提升后的顶点总数是 \(dV-\sum_P(e_P(F)-1)\)。局部标准形 \(t^e\) 说明提升的边和面在分歧点按有限个扇形拼接，给出 \(X\) 的有限胞腔分解；必要时再细分，不要求它一开始就是单纯剖分。于是
\[
 \chi(X)=d\chi(Y)-\sum_P(e_P(F)-1).
\]
将 \(\chi(X)=2-2g_X\)、\(\chi(Y)=2-2g_Y\) 代入并整理，即得公式。
\end{proof}

例如球面映射 \(z\mapsto z^d\) 的0与无穷远处重数均为 \(d\)，其余点重数为1；当 \(d>1\) 时恰有这两个分歧点，总分歧量 \(2d-2\) 抵消右侧的 \(-2d\)，留下 \(-2\)。若环面到球面的映射次数为2，每个分歧点只能有重数2，公式迫使分歧点恰有4个。若 \(X\rightarrow Y\) 无分歧，则 \(g_X-1=d(g_Y-1)\)；这个数值等式是必要条件，本身不构造覆盖。

\ProseParagraphBreak
除子语言把同一信息写成局部阶数的等式。对 \(Y\) 上除子 \(D=\sum_Q n_Q Q\)，定义
\[
 F^*D\coloneqq\sum_{P\in X}e_P(F)n_{F(P)}P,
 \quad R_F\coloneqq\sum_{P\in X}(e_P(F)-1)P.
\]
\(R_F\) 称为\term[fenqi-chuzi]{分歧除子}[ramification divisor]，纤维重数和公式给 \(\deg(F^*D)=d\deg D\)。由\cref{b3:prop:A-pullback-differential}，对 \(Y\) 上任意非零亚纯形式都有
\[
 (F^*\omega)=F^*((\omega))+R_F.
\]
这里尚未预设形式除子的次数。最后一节将在得到非常值亚纯函数之后，把此式与 Riemann--Hurwitz 接起来，算出该次数为 \(2g_X-2\)。

\section{多项式方程给出的曲面}
\label{b3:sec:A04}

\subsection{隐函数坐标与无穷远检查}
\label{b3:subsec:A0401}

多项式方程 \(P(x,y)=0\) 给出一维复几何对象，但“满足一个方程”还没有说明每一点附近都是一张光滑坐标图。若 \(P_y(x_0,y_0)\ne0\)，全纯隐函数定理给 \(y=h(x)\)，因而 \(x\) 是局部坐标；若 \(P_x\ne0\)，交换两变量即可。所用版本已在\subsecref{b3:subsec:220401}说明：一个非零复偏导对应可逆的二阶实矩阵，实隐函数定理先给实 \(C^1\) 图；对方程求导，再由 Cauchy--Riemann 方程得图函数全纯。这里需要在所讨论点检查非零偏导，不能只凭曲线的名字套用结论。

\ProseParagraphBreak
即使全部有限点都光滑，无穷远仍可能出问题。例如 \(y^2=x^5-1\) 在仿射平面没有两个偏导同时为0的点，普通射影闭包却由
\[
 Y^2Z^3=X^5-Z^5
\]
给出。在 \([0:1:0]\) 处，该齐次方程的三个一阶偏导都为0，原来的隐函数检查失效。下面直接构造合适的无穷远参数；我们要获得的是光滑紧模型，不能把普通平面射影闭包与它自动认作同一解析曲面。

\subsection{二重覆盖的光滑紧化}
\label{b3:subsec:A0402}

考虑 \(y^2=p(x)\)，其中 \(p\in\mathbb C[x]\) 无重根，\(N=\deg p\ge3\)。在 \(p(x)\ne0\) 的小圆盘中可以选两种平方根，它们给两张以 \(x\) 为坐标的互不相交的图。若 \(a\) 是简单根，写
\[
 p(x)=(x-a)q(x),\quad q(a)\ne0.
\]
选局部全纯根 \(v^2=q\)。曲线上令 \(t=y/v(x)\)，则
\[
 x=a+t^2,\quad y=t\,v(a+t^2).
\]
这张图同时包含原来的两层；投影到 \(x\) 的局部重数为2，曲面本身却在 \(t=0\) 光滑。

\ProseParagraphBreak
无穷远取何种坐标由 \(N\) 的奇偶决定。先设 \(N=2g+2\)。令 \(s=1/x\)、\(v=y s^{g+1}\)，方程成为
\[
 v^2=q_\infty(s),\quad
 q_\infty(s)=s^{2g+2}p(1/s).
\]
\(q_\infty\) 在0全纯且值为 \(p\) 的非零首项系数，所以有两种非零平方根。分别用 \(s\) 为参数，便补入两个无穷远点 \(\infty_+,\infty_-\)。目标球面在无穷远也用 \(s=1/x\)，所以两处局部重数均为1。

\ProseParagraphBreak
再设 \(N=2g+1\)。令 \(x=t^{-2}\)，并写
\[
 y=t^{-(2g+1)}u(t),\quad
 u(t)^2=t^{4g+2}p(t^{-2}).
\]
右侧在0全纯、非零而且为偶函数。选定 \(u(0)\) 后，平方根唯一性使 \(u\) 也是偶函数。于是 \(t\) 与 \(-t\) 对应相同 \(x\) 和相反 \(y\)，一张穿孔 \(t\) 圆盘恰好覆盖大 \(|x|\) 时的两层。因此只补一个无穷远点 \(\infty\)，目标坐标 \(1/x=t^2\) 表明其重数为2。若将平方根的两个符号各补一点，就会把同一张穿孔图重复计算。

\Needspace{6\baselineskip}
\begin{proposition}[二重覆盖的紧曲面模型]\label{b3:prop:A-hyperelliptic-model}
设 \(p\in\mathbb C[x]\) 无重根，\(\deg p=N\ge3\)，并令 \(g=\lfloor(N-1)/2\rfloor\)。仿射曲线 \(y^2=p(x)\) 由上述有限点坐标及无穷远坐标补成一个紧 Riemann 曲面 \(X\)。函数 \(x\) 延拓为次数2的全纯映射 \(X\rightarrow\widehat{\mathbb C\kern0pt}\)。其有限分歧点恰位于 \(p\) 的根上；\(N\) 为奇数时无穷远另有一个分歧点，\(N\) 为偶数时有两个无分歧的无穷远点。曲面 \(X\) 的亏格为 \(g\)。
\end{proposition}
\begin{proof}
有限部分在每点至少有一个非零偏导，给通常的子空间拓扑与上述全纯坐标。各无穷远穿孔图与大 \(|x|\) 部分通过所写的 \((x,y)\) 表达粘合，转换全纯且可逆；补入的点由所在图分离，偶次数的两点有互不相交的图。有限部分是 \(\mathbb C^2\) 的子空间，已有可数基，再加有限张无穷远图，得到 Hausdorff、第二可数的复曲面。

给定充分大的 \(R\)，有限部分中 \(|x|\le R\) 的点集在 \(\mathbb C^2\) 内闭而有界，因为 \(|y|^2=|p(x)|\)，所以紧。\(|x|\ge R\) 的部分连同补入点则由上述一个或两个闭坐标圆盘覆盖，也紧。两部分的并是 \(X\)，故 \(X\) 紧。

还需核对连通性。避开 \(p\) 的根后，任意有限点可沿平面路径的提升连到同一个基点上方的两点之一；绕某个简单根一周，平方根改变符号，连接这两个原像。根处的图和无穷远图都与该连通部分相交，故整个 \(X\) 连通。\(x\) 的一般纤维有两点，次数为2，局部图已经给出全部分歧。无论奇偶，总分歧点数都是 \(2g+2\)，各重数为2。Riemann--Hurwitz 因而给
\[
 2g_X-2=2(-2)+(2g+2)=2g-2,
\]
即 \(g_X=g\)。
\end{proof}

这些方程给\term[chao-tuoyuan-qumian]{超椭圆曲面}[hyperelliptic surface]的具体模型，也包含 \(g=1\) 的椭圆情形。例如次数3、4均给亏格1，次数5、6均给亏格2。次数为奇数时，有限根比偶数情形少一个，而无穷远分歧恰补上这个差额。

\subsection{全纯微分的具体计算}
\label{b3:subsec:A0403}

在刚构造的亏格 \(g\) 曲面上考虑
\[
 \omega_j=\frac{x^j\,\mathrm{d}x}{y},\quad 0\le j\le g-1.
\]
分母在根处为0，不能仅看这一表达就断言形式有极点。根 \(a\) 附近用 \(x=a+t^2\)、\(y=t\,v(a+t^2)\)，便有
\[
 \omega_j=\frac{2(a+t^2)^j}{v(a+t^2)}\,\mathrm{d}t,
\]
它在 \(t=0\) 全纯；其中 \(\omega_0\) 还处处非零。在有限非根点，以 \(x\) 为坐标时同样全纯。

\ProseParagraphBreak
无穷远的计算决定允许哪些幂次。偶次数时，\(x=s^{-1}\)、\(y=v(s)s^{-(g+1)}\)，所以
\[
 \omega_j=-\frac{s^{g-1-j}}{v(s)}\,\mathrm{d}s.
\]
奇次数时，\(x=t^{-2}\)、\(y=t^{-(2g+1)}u(t)\)，所以
\[
 \omega_j=-\frac{2t^{2g-2-2j}}{u(t)}\,\mathrm{d}t.
\]
给定范围内的指数均非负，这些形式因此全局全纯。若它们的一组线性组合为0，在任一有限非分歧图中除去 \(\mathrm{d}x/y\)，得到多项式 \(\sum c_jx^j\) 在开集上为0，所有系数只能为0。这里已直接得到 \(g\) 个独立形式；它们是否张成全部全纯形式，还要用最后一节的维数结论。

\ProseParagraphBreak
特别地，\(\mathrm{d}x/y\) 的全部零点都在无穷远，其除子为
\[
 \left(\frac{\mathrm{d}x}{y}\right)=
 \begin{cases}
 (g-1)(\infty_++\infty_-),&N=2g+2,\\
 (2g-2)\infty,&N=2g+1.
 \end{cases}
\]
两个式子的次数均为 \(2g-2\)。这个直接计算先在一族曲面上展示典范次数，稍后将证明任何紧 Riemann 曲面都有同样的规律。

\subsection{尖点、结点与归一化}
\label{b3:subsec:A0404}

多项式方程也可能把光滑参数压成奇点。尖点 \(y^2=x^3\) 有参数式 \((x,y)=(t^2,t^3)\)，只有一条局部分支。它与参数圆盘可以拓扑同胚，但原来的平面解析结构并不光滑：任一环境全纯函数 \(G(x,y)\) 的复合 \(G(t^2,t^3)-G(0,0)\)，若不恒为0，首个非零项阶数至少为2，因而不可能等于参数 \(t\)。恢复 \(t\) 是补入新的解析坐标，不是从原平面坐标中找一个非零导数。

\ProseParagraphBreak
结点 \(y^2=x^2(x+1)\) 则在原点有两条局部分支
\[
 y=x\sqrt{1+x},\quad y=-x\sqrt{1+x}.
\]
每一支单独可用 \(x\) 作光滑参数，但它们在平面中共享同一个原点。光滑模型要把这两个分支上的原点分开。将局部分支分离并恢复光滑参数的过程称为\term[guiyihua]{归一化}[normalization]。一般代数曲线的归一化存在性需要进一步的代数或解析构造。对带奇点曲线直接套用光滑平面曲线的亏格公式，会忽略归一化改变的局部信息。

\section{层与全局拼接}
\label{b3:sec:A05}

\subsection{限制、芽与相容截面}
\label{b3:subsec:A0501}

全纯函数是否全纯可以在每一点附近判定，若局部函数在重叠处相等，它们又能唯一拼接。层把这两件事与限制映射一并保存；重叠处的一致性正是把局部对象拼成全局对象的条件。

\Needspace{6\baselineskip}
\begin{definition}[层与截面]\label{b3:def:A-sheaf}
设 \(X\) 是拓扑空间。对每个开集 \(U\subset X\) 给一个集合或复向量空间 \(\mathcal F(U)\)，对每个包含关系 \(V\subset U\) 给限制映射 \(s\mapsto s|_V\)，并要求恒等限制不变、连续限制的复合等于直接限制；向量空间情形的限制映射还须线性。若对于 \(U\) 的任意开覆盖 \((U_i)\)，还满足以下两条性质，则称 \(\mathcal F\) 为 \(X\) 上的\term[ceng]{层}[sheaf]：两个对象若在所有 \(U_i\) 上限制相等，则在 \(U\) 上相等；一族 \(s_i\in\mathcal F(U_i)\) 若在所有重叠处限制相等，则存在 \(s\in\mathcal F(U)\) 限制为各 \(s_i\)。\(\mathcal F(U)\) 中的对象称为 \(U\) 上的\term[jiemian]{截面}[section]。
\end{definition}

在 Riemann 曲面上，\(\mathcal O(U)\) 表示 \(U\) 上全纯函数，\(\mathcal M(U)\) 表示亚纯函数，\(\Omega(U)\) 表示全纯1形式。它们的限制与拼接就是实际函数或形式的限制与拼接。给定除子 \(D\)，还可定义 \(\mathcal O(D)(U)\)：由满足 \(\operatorname{ord}_P f+D(P)\ge0\) 的亚纯函数组成，条件在每个 \(P\in U\) 检查。于是 \(\mathcal O(D)(X)=L(D)\)。

\ProseParagraphBreak
同一点附近的两个截面，若在更小的邻域相等，就代表同一个芽；所有这样的芽组成该点的\term[jing]{茎}[stalk]，记为 \(\mathcal F_P\)。在 \(P\) 的参数 \(t\) 中，若 \(D(P)=m\)，则
\[
 \mathcal O(D)_P=t^{-m}\mathcal O_P.
\]
意思是允许的局部亚纯函数都可以写为 \(t^{-m}\) 乘一个全纯函数。换参数只把局部生成元乘以一个不消失的全纯因子。最后用 \(L(D)\) 构造射影映射时，正是通过这种共同因子把极点消掉。

\ProseParagraphBreak
对复向量空间层 \(\mathcal F\) 及两个开集 \(U,V\)，层公理给出
\[
 0\longrightarrow\mathcal F(U\cup V)
 \longrightarrow\mathcal F(U)\oplus\mathcal F(V)
 \xrightarrow{\ \delta\ }\mathcal F(U\cap V),
 \quad \delta(s,t)=s|_{U\cap V}-t|_{U\cap V}
\]
在前两处精确：第一个映射单射，且 \(\delta\) 的核恰是可拼接的二元组。最后一个差值映射未必满射。局部数据能否被修正成相容的一族，问题就出在这里。

\subsection{局部差值与允许的修正}
\label{b3:subsec:A0502}

看 \(\mathbb C^*\) 上的形式 \(\mathrm{d}z/z\)。取两个割线域，分别用辐角区间 \((-\pi,\pi)\) 与 \((0,2\pi)\) 定义对数原函数 \(F_U,F_V\)。它们的交集有上下两个连通分支，上半平面有 \(F_U-F_V=0\)，下半平面则有 \(F_U-F_V=-2\pi\mathrm{i}\)。若只给 \(F_U,F_V\) 各加一个常数，两个差值不能同时变为0，所以这些局部原函数无法拼成全局原函数。

\ProseParagraphBreak
这里必须说清修正的类别。为了仍满足 \(\mathrm{d}F=\mathrm{d}z/z\)，只能加局部常数；因此上例的障碍属于局部常数形成的层。如果允许加任意全纯函数，导数也会改变，便不再解决原来的原函数问题。不能把这个例子直接说成全纯函数层 \(\mathcal O\) 的同一个障碍。

\ProseParagraphBreak
一般地，在固定开覆盖 \((U_i)\) 的两两交集上给差值 \(c_{ij}\)，要求 \(c_{ji}=-c_{ij}\)，并在三重交集满足
\[
 c_{ij}+c_{jk}=c_{ik}.
\]
这正是 \(c_{ij}=s_i-s_j\) 应有的相容关系，满足此关系的数据称为\term[yi-jie-shangbilian]{1阶上闭链}[1-cocycle]。若能找到允许的局部修正 \(b_i\)，使 \(c_{ij}=b_i-b_j\)，则 \(s_i-b_i\) 的差值全部为0，可以拼接。固定覆盖上相容差值模去这种修正的商，已记录一类拼接障碍；\term[ceng-shangtongdiao]{层上同调}[sheaf cohomology]把它与覆盖细化及更一般的分解统一起来。从固定覆盖的差值商到一般层上同调，还要处理覆盖的细化；随后使用的有限维性和对偶性是这一理论的进一步结果。

\subsection{主部问题与留数条件}
\label{b3:subsec:A0503}

紧曲面上的主部问题给出与 Riemann--Roch 直接相关的例子。在不同点 \(P_1,\ldots,P_s\) 选互不相交的小坐标盘 \(U_j\)，给定负次 Laurent 多项式
\[
 p_j(t_j)=\sum_{\nu=1}^{m_j}a_{j,\nu}t_j^{-\nu}.
\]
希望找到全局亚纯函数 \(f\)，使 \(f-p_j\) 在各 \(P_j\) 全纯，且在其余地方没有极点。记 \(U_0=X\setminus\{P_1,\ldots,P_s\}\)、\(p_0=0\)。在重叠处，\(p_i-p_j\) 全纯；若存在全纯修正 \(h_i\in\mathcal O(U_i)\) 使
\[
 p_i-p_j=h_i-h_j,
\]
则 \(p_i-h_i\) 拼成所求 \(f\)。反过来，若 \(f\) 存在，令 \(h_i=p_i-f\) 就得到这种修正。这里的障碍确实属于全纯函数的拼接问题。

\ProseParagraphBreak
全局留数和给出一族可直接检验的必要条件。对任一全纯1形式 \(\omega\)，由于全纯修正不改变乘积的留数，必须有
\[
 \sum_{j=1}^s\operatorname{Res}_{P_j}(p_j\omega)=0.
\]
因此留数和只取决于指定主部，而与怎样选择局部全纯余项无关。它把主部数据与全纯形式配成一个双线性表达；真正待回答的是，这些条件是否已经找出了所有障碍。

\ProseParagraphBreak
在球面上可以直接实现所有有限个主部：有限点的主部由 \((z-a)^{-\nu}\) 的线性组合给出，无穷远的主部由多项式给出，将它们相加即可。有限点的这些有理项在无穷远全纯，后加的多项式在各有限点也全纯，所以不会破坏先前的数据。一般紧曲面上没有这样的全局有理坐标。下一节的深层对偶输入将说明，刚才的留数条件仍是充分条件，并由此精确计数允许的自由度。

\section{Riemann--Roch 与曲面类型}
\label{b3:sec:A06}

\subsection{维数公式及其深层输入}
\label{b3:subsec:A0601}

前面已经知道允许极点的空间有限维，却不知道给定局部主部时究竟少了几个自由系数。\term[riemann-roch-dingli]{Riemann--Roch 定理}[Riemann--Roch theorem]把这个差额表示为一个微分形式空间的维数；它同时处理允许极点和要求零点两种条件。

\Needspace{6\baselineskip}
\begin{theorem}[Riemann--Roch，深层输入]\label{b3:thm:A-riemann-roch}
设 \(X\) 为拓扑亏格 \(g\) 的紧 Riemann 曲面，\(D\) 为 \(X\) 上除子，定义
\[
 L(D)=\{0\}\cup\{\,f\text{ 为非零亚纯函数}\mid(f)+D\ge0\,\},
 \quad \ell(D)=\dim_{\mathbb C}L(D).
\]
记 \(i(D)\) 为零形式及满足 \((\omega)\ge D\) 的非零亚纯1形式所成复向量空间的维数。则两个空间均有限维，并且
\[
 \ell(D)-i(D)=\deg D+1-g.
\]
在 \((\omega)\ge D\) 中，\(D\) 的正系数要求形式有相应零点，负系数则允许相应极点。
\end{theorem}
\begin{proofsketch}
先对有效除子 \(D=\sum_j m_jP_j\) 说明维数的来源。允许的主部空间由所有 \(p_j(t_j)=\sum_{\nu=1}^{m_j}a_{j,\nu}t_j^{-\nu}\) 组成，维数为 \(\deg D\)。全局实现的必要条件是对每个全纯形式 \(\omega\)，都有 \(\sum_j\operatorname{Res}_{P_j}(p_j\omega)=0\)。若 \(\omega=b_j(t_j)\,\mathrm{d}t_j\)，要让这一条件对所有主部恒为0，就必须使 \(b_j\) 的前 \(m_j\) 个 Taylor 系数都为0，亦即 \((\omega)\ge D\)。

深层输入把全纯形式空间的维数识别为拓扑亏格 \(g\)，并说明这些留数条件也充分。于是主部到全纯形式对偶空间的线性映射秩为 \(g-i(D)\)，可实现主部的空间维数为 \(\deg D-g+i(D)\)。同一主部的两个全局实现相差全纯函数，因紧性而相差常数；再加这一维，得到有效 \(D\) 的公式。

充分性不是局部 Laurent 展开本身的结论。层上同调把全纯差值的拼接障碍放入 \(H^1(X,\mathcal O)\)，\term[serre-duiou]{Serre 对偶}[Serre duality]说明，障碍为0当且仅当它与所有全纯形式的留数配对均为0。对于含负系数的除子，还需同时满足零点条件；使用 \(\mathcal O(D)\subset\mathcal O(D+P)\) 的局部系数商及其上同调精确序列，每加一个点，相应的 Euler 特征增加1，再逐点加减得到一般公式。不能把有效情形的自由主部计数直接用于负系数。

这一证明概要采用三项外部结果：有关上同调的有限维性及精确序列、上述留数对偶的非退化性、解析亏格与拓扑亏格的识别。最后一项依赖 \term[hodge-dingli]{Hodge 定理}[Hodge theorem]：亏格 \(g\) 曲面的一阶实上同调维数为 \(2g\)，由实调和1形式表示；复化后分为全纯形式及其共轭两个等维部分，故全纯形式维数为 \(g\)。相关上同调、对偶与 Riemann--Roch 论述见 Forster 的《Lectures on Riemann Surfaces》\cite[pp. 129--131, 135--139, 156--158]{Forster1981Lectures}。该书先以 \(\dim H^1(X,\mathcal O)\) 定义解析亏格，本附录用拓扑亏格，应用其维数公式时，须由 Hodge 定理把解析亏格识别为这里的拓扑亏格。
\end{proofsketch}

公式中的修正项有具体含义。极点条件较少时，全纯形式会对主部施加限制；极点条件足够多以后，不再有非零形式同时在指定点消失那么多阶，\(i(D)\) 便为0，维数随 \(\deg D\) 以斜率1增长。下一条结果给出这一转折的准确阈值。

\subsection{亚纯函数、典范除子与高次数情形}
\label{b3:subsec:A0602}

\Needspace{6\baselineskip}
\begin{proposition}[函数存在与典范次数]\label{b3:prop:A-riemann-roch-consequences}
设 \(X\) 为亏格 \(g\) 的紧 Riemann 曲面。任给 \(P\in X\)，存在非常值亚纯函数，仅在 \(P\) 可能有极点，极点阶数至多 \(g+1\)。每个非零亚纯1形式的除子次数均为 \(2g-2\)，全纯1形式空间的复维数为 \(g\)。对任意除子 \(D\)，若 \(\deg D>2g-2\)，则满足 \((\omega)\ge D\) 的亚纯形式只能为0，并有
\[
 \ell(D)=\deg D+1-g,
\]
其中 \(\ell(D)\) 是由零函数及满足 \((f)+D\ge0\) 的非零亚纯函数所成空间的维数。
\end{proposition}
\begin{proof}
取 \(D=(g+1)P\)，Riemann--Roch 给 \(\ell(D)=2+i(D)\ge2\)，而常数只占一维，故有非常值 \(f\)。若它在 \(P\) 也无极点，就会全局全纯而为常数，因而 \(P\) 确为其极点。

把 \(f\) 看作次数 \(d\) 的映射 \(X\rightarrow\widehat{\mathbb C\kern0pt}\)。在有限值点，\(\mathrm{d}f\) 的阶数为 \(e_P(f)-1\)；在极点，写 \(f=t^{-e_P}u(t)\)，其中 \(u(0)\ne0\)，则 \(\mathrm{d}f\) 的阶数为 \(-e_P-1=(e_P-1)-2e_P\)。极点的局部重数之和为 \(d\)，所以
\[
 \deg(\mathrm{d}f)=\sum_P(e_P(f)-1)-2d=2g-2,
\]
末步用 Riemann--Hurwitz。任一其他非零亚纯形式与 \(\mathrm{d}f\) 的商都是全局亚纯函数，因为分子分母的坐标转换因子相消；其主除子次数为0，所以所有非零形式都有同一次数。

若 \((\omega)\ge D\) 且 \(\deg D>2g-2\)，有效除子 \((\omega)-D\) 的次数会为负，不可能。因此 \(i(D)=0\)，由 Riemann--Roch 得高次数公式。最后令 \(D=0\)，由 \(\ell(0)=1\)，同一公式给 \(i(0)=g\)，也就是全纯1形式空间的维数。
\end{proof}

固定一个非零亚纯形式 \(\omega_0\)，称 \(K_X=(\omega_0)\) 为一个\term[dianfan-chuzi]{典范除子}[canonical divisor]。换形式只把它加上一个主除子，故其线性等价类不变。映射 \(\omega\mapsto\omega/\omega_0\) 给出
\[
 i(D)=\ell(K_X-D),\quad
 \ell(D)-\ell(K_X-D)=\deg D+1-g.
\]
确实，\((\omega)\ge D\) 等价于 \((\omega/\omega_0)+K_X-D\ge0\)。这个常用写法把两个空间都用函数表示，但形式在留数配对中的作用仍不可省去。

\ProseParagraphBreak
\secref{b3:sec:A04}已构造的 \(g\) 个独立全纯形式，现在因维数恰为 \(g\) 而成为一组基。亏格0时取 \(D=P\)，有 \(\ell(P)=2\)，因而存在仅有一个简单极点的函数；它到球面的次数为1，所以 \(X\) 双全纯于球面。亏格1时，任一非零全纯形式的有效零点除子次数为0，因此处处不为0，典范除子线性等价于0；任意正次数除子都满足 \(\ell(D)=\deg D\)。

\subsection{局部插值与射影嵌入}
\label{b3:subsec:A0603}

维数公式还能够回答插值问题：除了允许指定极点，能否让局部正规化后的函数在若干点具有任意指定的有限个 Taylor 系数？若把增加的条件写成另一个有效除子，只需比较两个空间的维数。

\Needspace{6\baselineskip}
\begin{proposition}[有限阶局部插值]\label{b3:prop:A-jet-interpolation}
设 \(X\) 是亏格 \(g\) 的紧 Riemann 曲面，\(D\) 为除子，\(E=\sum_{j=1}^s m_jP_j\) 为有效除子，其中 \(P_j\) 两两不同，\(m_j\ge1\)。在各点取参数 \(t_j(P_j)=0\)。若
\[
 \deg D>2g-2+\deg E,
\]
则对每组给定复系数 \(a_{j,0},\ldots,a_{j,m_j-1}\)，均存在 \(f\in L(D)\)，使
\[
 t_j^{D(P_j)}f
 =\sum_{\nu=0}^{m_j-1}a_{j,\nu}t_j^\nu+O(t_j^{m_j})
\]
在各 \(P_j\) 成立。所有实现的差恰属于 \(L(D-E)\)。
\end{proposition}
\begin{proof}
取 \(f\in L(D)\) 的这些有限 Taylor 系数，得到线性映射
\[
 J:L(D)\longrightarrow\bigoplus_{j=1}^s\mathbb C^{m_j}.
\]
系数全为0等价于各 \(P_j\) 比 \(D\) 原来要求的阶数再多消失 \(m_j\) 阶，所以 \(\ker J=L(D-E)\)。假设使 \(D\) 与 \(D-E\) 的次数都大于 \(2g-2\)，高次数维数公式给
\[
 \dim\operatorname{im}J=\ell(D)-\ell(D-E)=\deg E.
\]
这恰为目标空间维数，故 \(J\) 满射。核的描述也给出实现的全部不唯一性。
\end{proof}

当插值点不在 \(D\) 的支集上时，正规化因子为1，上式就是普通函数值和导数的插值；当它位于允许极点处时，指定的是消去极点后的首批系数。后者也适用于 \(D(P_j)<0\)，因为原来的零点条件保证除去相应幂次以后仍全纯。

\ProseParagraphBreak
复射影空间 \(\mathbb{CP}^N\) 的点是非零向量 \((a_0,\ldots,a_N)\) 模去共同非零倍数后的类，记为 \([a_0:\cdots:a_N]\)。在 \(a_k\ne0\) 的图中，用比值 \(a_j/a_k\) 作仿射坐标。这种共同倍数不改变点的性质，正好消除 \(\mathcal O(D)\) 换局部生成元时的差别。要得到嵌入，还需分别排除公共零点、不能分离两个点以及不能分离切向这三种障碍。

\Needspace{6\baselineskip}
\begin{proposition}[由除子构造射影嵌入]\label{b3:prop:A-projective-embedding}
设 \(X\) 是亏格 \(g\) 的紧 Riemann 曲面，\(D\) 是 \(X\) 上除子，\(\deg D\ge2g+1\)。取由零函数及满足 \((f)+D\ge0\) 的非零亚纯函数所成空间 \(L(D)\) 的一组基 \(f_0,\ldots,f_N\)。在 \(P\) 的局部参数 \(t(P)=0\) 中，将各坐标同时乘以 \(t^{D(P)}\)，则
\[
 \Phi=[t^{D(P)}f_0:\cdots:t^{D(P)}f_N]
\]
全纯且没有公共零点。这些表达给出全局全纯嵌入 \(\Phi:X\rightarrow\mathbb{CP}^N\)。
\end{proposition}
\begin{proof}
设 \(m=D(P)\)。\(a_j=t^m f_j\) 都全纯，包括 \(m<0\) 的情况。对 \(E=P\) 用插值命题，局部取值泛函 \(E_P(f)=(t^m f)(P)\) 满射，故至少一个 \(a_j(P)\ne0\)。换参数时，全部 \(a_j\) 乘同一个处处非零的全纯因子，射影值不变；在不含除子其他支点的小图上，所得表达因此相容，定义全局全纯映射。

若 \(P\ne Q\)，对 \(E=P+Q\) 用插值，可选 \(f\in L(D)\) 使 \(E_P(f)=0\)、\(E_Q(f)=1\)。若 \(\Phi(P)=\Phi(Q)\)，两点所有基函数的正规化取值向量成非零比例，所以每个线性组合的取值也成同一比例，与所选 \(f\) 矛盾。因此 \(\Phi\) 单射。

固定 \(P\)，选一个 \(f_0\) 使 \(a_0(P)\ne0\)，必要时换基；换基只在目标射影空间复合一个可逆线性变换。对 \(E=2P\) 插值，可选 \(f\) 使 \(t^m f=t+O(t^2)\)。于是仿射图中的比值 \(f/f_0=(t^m f)/a_0\) 在 \(P\) 的导数不为0。它是基函数比值的线性组合，故 \(\Phi\) 的导数不为0。最后，紧空间到 Hausdorff 空间的连续单射是到其像的同胚；有一个非零坐标导数，又使像局部成为一个全纯图。因此 \(\Phi\) 是全纯嵌入。
\end{proof}

球面上取 \(D=2[\infty]\)，基 \(1,z,z^2\) 给映射 \(z\mapsto[1:z:z^2]\)。在无穷远令 \(t=1/z\)，同时乘以 \(t^2\)，便得到 \([t^2:t:1]\)，在 \(t=0\) 也有定义。像满足 \(X_0X_2=X_1^2\)，是一条光滑圆锥曲线。这个算例展示了为什么不能简单把 \(f_j(P)\) 直接代入极点，而必须先消去共同的局部阶数。

\ProseParagraphBreak
一般嵌入所得的是紧射影解析曲线。从这种解析描述进一步得到齐次多项式方程，是紧射影解析子集的代数性定理所解决的问题。反方向，无奇点的复射影代数曲线由隐函数坐标给出紧 Riemann 曲面；奇点必须先作归一化。Miranda 的《Algebraic Curves and Riemann Surfaces》\cite{Miranda1995Algebraic}可作为代数曲线、局部环和光滑模型的进阶阅读。

\subsection{亏格一的周期表示与曲面类型}
\label{b3:subsec:A0604}

亏格0已由一次亚纯函数识别为球面。亏格1时，处处无零的全纯形式提供局部坐标：积分后每张图的转换只差一个常数。剩下的问题是这些常数是否组成二维格点，以及积分映射会不会多次覆盖同一个环面。周期的面积恒等式同时解决这两点。

\Needspace{6\baselineskip}
\begin{proposition}[亏格一曲面的周期格点]\label{b3:prop:A-genus-one-torus}
设 \(X\) 为亏格1的紧 Riemann 曲面。则存在两个实线性无关的复数 \(A,B\)，使 \(X\) 双全纯于复环面 \(\mathbb C/\Lambda\)，其中 \(\Lambda=\mathbb ZA+\mathbb ZB\)。
\end{proposition}
\begin{proof}
由\cref{b3:prop:A-riemann-roch-consequences}，取非零全纯形式 \(\omega\)，其零点除子次数为0，所以处处无零。按拓扑环面的标准多边形表示，取正向生成曲线 \(\alpha,\beta\)，沿它们切开成为边界词为 \(\alpha\beta\alpha^{-1}\beta^{-1}\) 的四边形。令
\[
 A=\int_\alpha\omega,\quad B=\int_\beta\omega.
\]
切开内部单连通，\(\omega\) 有原函数 \(F\)；它沿各边延伸到边界，配对边上的原函数值相差常数，微分相同。曲线可取分段光滑，角点不影响积分。

将 \(F\) 加常数，使边界四个连续顶点的值为 \(0,A,A+B,B\)。由 Stokes 公式及
\(\mathrm{d}(\overline F\,\mathrm{d}F)=-\omega\wedge\overline\omega\)，有
\[
 \int_X\frac{\mathrm{i}}2\omega\wedge\overline\omega
 =\frac1{2\mathrm{i}}\int_{\partial}\overline F\,\mathrm{d}F.
\]
两条 \(\alpha\) 边反向配对，第二条上的 \(F\) 比第一条多 \(B\)，合计积分为 \(-\overline B A\)；两条 \(\beta\) 边合计为 \(\overline A B\)。因此
\[
 \int_X\frac{\mathrm{i}}2\omega\wedge\overline\omega
 =\operatorname{Im}(\overline A B)>0.
\]
最后的不等式来自局部表达 \((\mathrm{i}/2)\omega\wedge\overline\omega=|a(z)|^2\,\mathrm{d}x\,\mathrm{d}y\)，其中 \(\omega=a(z)\,\mathrm{d}z\) 处处非零。这证明 \(A,B\) 实线性无关，所以它们生成离散格点群 \(\Lambda\)。

固定 \(P_0\in X\)，定义
\[
 \Psi(P)=\int_{P_0}^P\omega\pmod\Lambda.
\]
任意闭路的同调类都是 \(\alpha,\beta\) 的整数线性组合；\(\omega\) 局部闭，由 Stokes 公式，其积分只依赖同调类，所以每个周期都属于 \(\Lambda\)。这保证 \(\Psi\) 与路径选择无关。局部用原函数表示时，\(\Psi\) 的导数为非零的 \(\omega\)，故它是从紧曲面到连通环面的有限无分歧覆盖，记次数为 \(d\)。

目标环面的平坦面积形式拉回为 \((\mathrm{i}/2)\omega\wedge\overline\omega\)，一个基本平行四边形的面积为 \(\operatorname{Im}(\overline A B)\)。将目标分成覆盖图，逐图换元，每片恰有 \(d\) 个原像，得到
\[
 \int_X\frac{\mathrm{i}}2\omega\wedge\overline\omega
 =d\operatorname{Im}(\overline A B).
\]
与前面的面积恒等式比较，因为该面积严格为正，所以 \(d=1\)，\(\Psi\) 双全纯。
\end{proof}

这项周期构造包含两个关键环节：正面积保证两条周期实线性无关，从而生成离散格；覆盖的面积计数则确定覆盖只有一层，使积分映射成为双全纯同构。

\ProseParagraphBreak
\cref{b3:thm:12-uniformization}从万有覆盖给出另一种描述：单连通覆盖模型是球面、平面或圆盘。结合\cref{b3:prop:12-compact-genus-types}，紧曲面的球面型就是亏格0的球面，平面型就是亏格1的二维格点商，其余亏格至少2的紧曲面属于圆盘型。一般统一化仍为前章标明的深层输入，本附录没有通过 Riemann--Roch 重新证明它。局部标准形描述点附近的映射，Riemann--Hurwitz 计算分歧的拓扑代价，Riemann--Roch 则计算允许局部数据时的全局自由度，三者的职责各不相同。

\begin{chaptersummary}
复结构由双全纯坐标转换确定，微分形式则按导数规律转换。局部重数同时控制纤维计数、拉回阶数与留数；紧性使局部原像层拼成有限次数，继而给出 Riemann--Hurwitz 公式。除子把零极点条件写成有限维线性问题，主部与全纯形式的留数配对揭示拼接障碍。Riemann--Roch 将拼接资料汇成维数公式，由此得到亚纯函数存在、典范次数、局部插值、射影嵌入与亏格一的周期格点表示。使用代数方程时，有限点、无穷远及奇点分支必须逐处检查，光滑模型不能被未经核验的射影闭包代替。
\end{chaptersummary}

\BookExercises

以下八题均为自拟。除直接计算外，题组依次训练主部实现、商曲面构造、典范映射、极点间隙、嵌入阈值、周期障碍与奇点归一化；所需主线结论均已在正文建立。

% 自拟：将局部拉回公式用于有分歧的有理映射，区分形式极点阶数与留数的变化。
\Needspace{6\baselineskip}
\begin{exercise}\label{b3:ex:A-1}
在球面上令 \(R(z)=z+1/z\)，并在目标球面取 \(\omega=\mathrm{d}w/(w^2-4)\)。求 \(R\) 的次数、全部分歧点及分歧值，计算 \(R^*\omega\) 的留数和除子，逐点验证
\[
 (R^*\omega)=R^*((\omega))+R_R.
\]
必须检查 \(z=0\) 与无穷远，不能只在有限且分母非零处约分。
\end{exercise}
\begin{solution}
\(R\) 在0和无穷远各有一个简单极点，故次数为2。有限非零处
\[
 R'(z)=1-z^{-2},\quad
 R(z)-2=\frac{(z-1)^2}{z},\quad
 R(z)+2=\frac{(z+1)^2}{z}.
\]
所以 \(z=1,-1\) 的局部重数均为2，分歧值分别为 \(2,-2\)。在0用目标坐标 \(1/R=z/(1+z^2)\)，导数为1；在无穷远用 \(t=1/z\)，同样有 \(1/R=t/(1+t^2)\)，重数为1。因此无其他分歧点。

目标形式在 \(w=2,-2\) 分别有留数 \(1/4,-1/4\)。无穷远用 \(s=1/w\)，有 \(\omega=-\mathrm{d}s/(1-4s^2)\)，所以此处全纯且非零，\((\omega)=-[2]-[-2]\)。在可直接代入的开集上计算
\[
 R^*\omega
 =\frac{1-z^{-2}}{(z+z^{-1})^2-4}\,\mathrm{d}z
 =\frac{\mathrm{d}z}{z^2-1}.
\]
两侧为同一亚纯形式的表达，故等式延拓到所有点。后式在 \(1,-1\) 的留数为 \(1/2,-1/2\)，恰为目标留数乘局部重数2；在0全纯且系数为 \(-1\)。无穷远改用 \(t=1/z\) 后变成 \(-\mathrm{d}t/(1-t^2)\)，也全纯且非零。因此 \((R^*\omega)=-[1]-[-1]\)。另一方面，\(R^*((\omega))=-2[1]-2[-1]\)，\(R_R=[1]+[-1]\)，相加正得该除子。
\end{solution}

% 自拟：在亏格一情形用维数与留数的双重计数证明指定简单主部的充分性。
\Needspace{6\baselineskip}
\begin{exercise}\label{b3:ex:A-2}
设 \(X\) 是亏格1的紧 Riemann 曲面，\(\omega\ne0\) 为全纯1形式，\(P_1,\ldots,P_n\) 为不同点，\(n\ge1\)。在各点取局部参数 \(t_j(P_j)=0\)，使 \(\omega=\mathrm{d}t_j\)。证明：存在仅在这些点可能有极点、且主部依次为 \(c_j/t_j\) 的全局亚纯函数，当且仅当 \(\sum_j c_j=0\)。有实现时说明其唯一性范围，并判断哪些点确为极点。
\end{exercise}
\begin{solution}
\(\omega\) 处处无零，局部积分的导数非零，所以所要求的参数存在。若 \(f\) 实现这些主部，则 \(f\omega\) 在 \(P_j\) 的留数为 \(c_j\)；全局留数和给 \(\sum_j c_j=0\)。

令 \(D=P_1+\cdots+P_n\)。亏格1且 \(\deg D=n>0\)，故 \(\ell(D)=n\)。把函数映到各简单主部系数，得到线性映射 \(T:L(D)\rightarrow\mathbb C^n\)。其核由没有极点的函数组成，正是常数，维数为1。因此 \(\dim\operatorname{im}T=n-1\)。必要性已证明像包含于超平面 \(\{(c_j):\sum_j c_j=0\}\)，该超平面同样为 \(n-1\) 维，所以两者相等，这就给充分性，包括 \(n=1\) 时只有零主部的情况。

两个实现之差没有任何极点，必为常数，故实现唯一到加常数。\(c_j\ne0\) 时 \(P_j\) 是简单极点，\(c_j=0\) 时因没有更高负次项，该点可去。
\end{solution}

% 自拟：先构造二阶自同构的商坐标，再用 Riemann--Hurwitz 计算固定点。
\Needspace{6\baselineskip}
\begin{exercise}\label{b3:ex:A-3}
设 \(X\) 为亏格 \(g_X\) 的紧 Riemann 曲面，\(\sigma:X\rightarrow X\) 为满足 \(\sigma^2=\operatorname{id}\)、\(\sigma\ne\operatorname{id}\) 的全纯自同构。构造商空间 \(Y=X/\langle\sigma\rangle\) 的复结构，使商映射次数为2。若固定点数为 \(r\)，证明
\[
 r=2g_X+2-4g_Y.
\]
由此说明 \(r\) 必为偶数，\(r=0\) 时 \(g_X\) 为奇数，以及 \(r=2g_X+2\) 当且仅当 \(Y\) 双全纯于球面。
\end{exercise}
\begin{solution}
非固定点可取小坐标盘 \(U\) 使 \(U\cap\sigma(U)=\varnothing\)，于是商映射在 \(U\) 上给商的一个普通坐标。若 \(P\) 固定，在 \(z(P)=0\) 的局部表达中，\(\sigma^2=\operatorname{id}\) 先给 \(\sigma'(0)^2=1\)。若 \(\sigma'(0)=1\)，而 \(\sigma(z)-z\) 首个非零项为 \(a_mz^m\)，\(m\ge2\)，则复合后 \(\sigma^2(z)=z+2a_mz^m+O(z^{m+1})\)，矛盾。若根本没有非零项，恒等定理使 \(\sigma\) 全局为恒等，也矛盾。所以 \(\sigma'(0)=-1\)。

在缩小的 \(\sigma\) 不变邻域内令
\[
 t=\frac{z-z\circ\sigma}{2}.
\]
则 \(t'(P)=1\)、\(t\circ\sigma=-t\)。商坐标可取 \(u=t^2\)，其纤维恰为 \(\sigma\) 的轨道，商映射在此重数为2。不同商坐标的转换全纯：在自由轨道处来自普通图的全纯转换；固定点附近，不变全纯函数的 \(t\) 展开只有偶次项，因而是 \(u=t^2\) 的全纯函数，坐标转换可逆。

还要检查商的拓扑性质。有限轨道的不同点集可由不交邻域分开，再取有限次平移的交并给不交不变邻域，故商 Hausdorff。商映射开，因为一个开集的饱和为 \(U\cup\sigma(U)\)；原空间可数基的像给商的可数基。紧性与连通性由连续满射保持，因此 \(Y\) 是紧 Riemann 曲面。固定点局部是 \(\sigma-\operatorname{id}\) 的孤立零点，紧性使其有限。自由轨道有两点，商映射次数为2，其全部分歧点正是固定点。

Riemann--Hurwitz 给 \(2g_X-2=2(2g_Y-2)+r\)，整理得所求公式。因此 \(r\) 为偶数，且 \(0\le r\le2g_X+2\)；若 \(r=0\)，则 \(g_X=2g_Y-1\) 为奇数。等号 \(r=2g_X+2\) 等价于 \(g_Y=0\)，再由正文的亏格0分类，等价于商曲面为球面。
\end{solution}

% 自拟：利用已构造的微分基研究典范映射，训练无公共零点与分离点的区别。
\Needspace{6\baselineskip}
\begin{exercise}\label{b3:ex:A-4}
设 \(X\) 是 \(y^2=p(x)\) 的光滑紧模型，其中 \(p\) 无重根，次数为 \(2g+1\) 或 \(2g+2\)，\(g\ge2\)。用全纯形式基 \(\omega_j=x^j\,\mathrm{d}x/y\)、\(0\le j\le g-1\)，定义射影映射 \(\kappa=[\omega_0:\cdots:\omega_{g-1}]\)，其意义是在任一局部坐标中取形式系数的射影比。证明它处处有定义，并分解为二重投影 \(x:X\rightarrow\widehat{\mathbb C\kern0pt}\) 后接映射
\[
 z\longmapsto[1:z:\cdots:z^{g-1}].
\]
求无穷远的射影值，说明 \(\kappa\) 为何不是嵌入，并说明 \(g=2\) 的特殊情形。
\end{exercise}
\begin{solution}
换局部坐标时各形式的系数都乘同一个非零转换导数，所以射影比不变。有限点上 \(\omega_0\) 无零点，包括 \(y=0\) 的分歧点；无穷远的计算则显示 \(\omega_{g-1}\) 无零点。因此没有所有形式同时消失的点，\(\kappa\) 处处有定义。

有限点除去共同因子 \(\mathrm{d}x/y\)，得 \(\kappa=[1:x:\cdots:x^{g-1}]\)。偶次数的两个无穷远点处，\(\omega_j\) 的阶数为 \(g-1-j\)；奇次数的单个无穷远点处为 \(2g-2-2j\)。因此在那里都只有最后一个坐标可能非零，值均为 \([0:\cdots:0:1]\)。球面映射在无穷远用 \(s=1/z\) 写成 \([s^{g-1}:s^{g-2}:\cdots:1]\)，正好给出同一值，故分解全局成立。

后一个球面映射是嵌入：有限处可由第二坐标与第一坐标之比恢复 \(z\)，无穷远附近可由倒数第二坐标与最后坐标之比恢复 \(s\)。但前一个映射在非分歧值上把 \((x,y)\) 与 \((x,-y)\) 送到同一点，\(\kappa\) 因此不能分离这两个点，不是嵌入。\(g=2\) 时后一个映射就是球面的通常射影表示，\(\kappa\) 本身为次数2的映射。这说明全纯形式给出的典范坐标即使没有公共零点，也未必嵌入曲面。
\end{solution}

% 自拟：由单点增维和高次数公式建立极点间隙，再用于奇次超椭圆模型。
\Needspace{6\baselineskip}
\begin{exercise}\label{b3:ex:A-5}
设 \(X\) 为亏格 \(g\ge1\) 的紧 Riemann 曲面，\(P\in X\)。若不存在仅在 \(P\) 有极点且阶数恰为正整数 \(m\) 的亚纯函数，就称 \(m\) 为 \(P\) 的一个\term[jianxi]{间隙}[gap]。证明间隙恰有 \(g\) 个，并全部位于 \(1,\ldots,2g-1\)；证明正的非间隙在加法下封闭。对于奇次模型 \(y^2=p(x)\)、\(\deg p=2g+1\)，求其唯一无穷远点的全部间隙。
\end{exercise}
\begin{solution}
函数有指定的精确极点阶数 \(m\)，等价于 \(L(mP)\setminus L((m-1)P)\) 非空。由正文的单点系数映射，增量
\[
 \delta_m=\ell(mP)-\ell((m-1)P)
\]
只能为0或1，且间隙恰对应 \(\delta_m=0\)。当 \(m\ge2g-1\) 时，高次数公式给 \(\ell(mP)=m+1-g\)。特别地，\(\ell((2g-1)P)=g\)，而 \(\ell(0)=1\)，所以前 \(2g-1\) 个增量的和为 \(g-1\)，其中恰有 \((2g-1)-(g-1)=g\) 个为0。对 \(m\ge2g\)，相邻两项都可用高次数公式，增量为1，没有新的间隙。

若 \(a,b\) 是正非间隙，取分别仅在 \(P\) 有精确 \(a,b\) 阶极点的函数 \(f,h\)，则 \(fh\) 仅在 \(P\) 有精确 \(a+b\) 阶极点，故非间隙对加法封闭。

奇次超椭圆模型在无穷远有 \(x=t^{-2}\)、\(y=t^{-(2g+1)}u(t)\)，\(u(0)\ne0\)。所以2和 \(2g+1\) 都是非间隙，\(x^k\) 给所有正偶数，\(yx^k\) 给所有不小于 \(2g+1\) 的奇数。因此间隙只能位于 \(1,3,\ldots,2g-1\)。这组数恰有 \(g\) 个，而总间隙数已证为 \(g\)，故它们全部且恰好是间隙。
\end{solution}

% 自拟：检查射影嵌入阈值少一阶时的精确障碍，使用次数零除子的判别。
\Needspace{6\baselineskip}
\begin{exercise}\label{b3:ex:A-6}
设 \(X\) 为亏格 \(g\ge1\) 的紧 Riemann 曲面，\(K_X\) 为典范除子，\(D\) 为次数 \(2g\) 的除子。证明 \(L(D)\) 的任一基用正文方法给出一个处处有定义的射影映射。进一步证明：对于不同点 \(P,Q\)，该映射不能分离 \(P,Q\) 当且仅当 \(D\sim K_X+P+Q\)；它在 \(P\) 的微分为0当且仅当 \(D\sim K_X+2P\)。据此解释 \(2g+1\) 的嵌入条件比 \(2g\) 多排除了什么。
\end{exercise}
\begin{solution}
高次数公式给 \(\ell(D)=g+1\)、\(\ell(D-P)=g\)，所以每点正规化取值泛函均非零，基函数没有公共零点，射影映射处处有定义。

对不同 \(P,Q\)，两点取值映射的核为 \(L(D-P-Q)\)。Riemann--Roch 给
\[
 \ell(D-P-Q)=g-1+\ell(K_X-D+P+Q).
\]
右侧最后一个除子次数为0。由正文的次数0判别，其空间维数仅可能为0或1，且为1当且仅当该除子为主除子，也就是 \(D\sim K_X+P+Q\)。因此两点取值映射的秩分别为2或1。两个单点取值泛函都非零，秩为1恰好表示二者成非零比例，亦即两点的射影坐标相同，得到第一条等价关系。

在 \(P\) 取参数 \(t\)，把 \(t^{D(P)}f\) 的常数项和一次项作为二元数据。这个映射的核为 \(L(D-2P)\)，同样有
\[
 \ell(D-2P)=g-1+\ell(K_X-D+2P).
\]
故两项数据的秩为1当且仅当 \(D\sim K_X+2P\)。选基函数 \(f_0\) 的正规化值在 \(P\) 非零。秩为1等价于每个正规化函数的一次项都由其常数项通过同一比例决定；用商的求导公式，这恰等价于全部 \(f/f_0\) 在 \(P\) 的导数为0，即射影映射微分为0。秩为2则存在常数项为0、一次项非零的 \(f\)，相应比值导数不为0。

次数 \(2g\) 已能排除公共零点，但可能留下上述次数0除子为主除子的特殊障碍。若次数升为 \(2g+1\)，两点或二阶局部数据所对应的修正除子次数都变为负，修正空间必为0，两个障碍同时消失。
\end{solution}

% 自拟：球面给充分性，正亏格曲面构造只有一个极点而不全局恰当的形式。
\Needspace{6\baselineskip}
\begin{exercise}\label{b3:ex:A-7}
证明球面上的亚纯1形式具有全局亚纯原函数，当且仅当它在每点的留数均为0。再证明：对任意亏格 \(g\ge1\) 的紧 Riemann 曲面 \(X\) 和任意点 \(P\)，都存在一个恰在 \(P\) 有极点、所有留数为0、却没有全局亚纯原函数的亚纯1形式。说明两种情形的差异所在。
\end{exercise}
\begin{solution}
若 \(\omega=\mathrm{d}F\)，对 \(F\) 的 Laurent 展开逐项微分可见 \(z^{-1}\,\mathrm{d}z\) 项的系数为0，故留数全为0。反过来，球面上写 \(\omega=a(z)\,\mathrm{d}z\)，系数 \(a\) 为有理函数。将 \(a\) 作部分分式分解，留数为0使全部一阶主部项消失，故
\[
 a(z)=q(z)+\sum_{j}\sum_{k=2}^{m_j}\frac{c_{j,k}}{(z-a_j)^k},
\]
其中 \(q\) 是多项式。逐项积分，多项式有多项式原函数，余项的原函数为 \(c_{j,k}(z-a_j)^{1-k}/(1-k)\)。所得 \(F\) 为有理函数，在有限开集满足 \(\mathrm{d}F=\omega\)，故作为球面上的亚纯形式也处处相等。

现在设 \(g\ge1\)。由正文取一个非零全纯形式 \(\eta\)，以及仅在指定点 \(P\) 有极点的非常值亚纯函数 \(f\)。令 \(\omega=\mathrm{d}f+\eta\)。若 \(f\) 在 \(P\) 有 \(m\) 阶极点，\(\mathrm{d}f\) 在那里有 \(m+1\) 阶极点，不能被全纯的 \(\eta\) 抵消；其余点无极点。两项的留数都为0，因此 \(\omega\) 的留数全为0。

若有全局亚纯函数 \(F\) 满足 \(\mathrm{d}F=\omega\)，则 \(\mathrm{d}(F-f)=\eta\) 全纯。亚纯函数的极点微分后仍是极点，所以 \(F-f\) 不可能有极点；紧性使它为常数，进而 \(\eta=0\)，矛盾。球面没有非零全纯1形式，部分分式中的留数条件足以构造原函数；正亏格曲面则还有全局周期障碍，本构造把该障碍保留在 \(\eta\) 中。
\end{solution}

% 自拟：保留结点参数化，增加归一化上微分的极点与相反留数计算。
\Needspace{6\baselineskip}
\begin{exercise}\label{b3:ex:A-8}
用斜率 \(t=y/x\) 给结点曲线 \(y^2=x^2(x+1)\) 构造有理参数，并延伸为从球面到其射影闭包的映射。找出结点和无穷远点的全部原像，说明光滑模型为何为亏格0。最后计算 \(\mathrm{d}x/y\) 拉回到参数球面后的零极点及留数，解释为什么原平面上的这个表达不能直接宣布为光滑模型上的全纯形式。
\end{exercise}
\begin{solution}
在 \(x\ne0\) 处有 \(t^2=x+1\)，故 \(x=t^2-1\)、\(y=t(t^2-1)\)。结点 \((0,0)\) 的原像为 \(t=1,-1\)，两点分别对应两条局部分支。射影方程为 \(Y^2Z=X^2(X+Z)\)。写 \(t=T/U\)，参数化成为
\[
 [T:U]\longmapsto
 [\,U(T^2-U^2):T(T^2-U^2):U^3\,].
\]
若 \(U\ne0\)，第三个坐标非零；若 \(U=0\)，则 \(T\ne0\)，第二个坐标为 \(T^3\ne0\)。所以这些齐次多项式无公共零点，映射全局有定义。射影曲线上 \(Z=0\) 迫使 \(X=0\)，故唯一无穷远点为 \([0:1:0]\)，其唯一原像是 \(U=0\)。在 \(Y\ne0\) 的图中，沿参数化有 \(X/Y=U/T\)，因此无穷远参数仍给非退化局部坐标；有限非结点处则由 \(t=y/x\) 反解。参数球面分开结点的两支并补上无穷远，得到光滑模型，亏格为0。光滑平面三次曲线的亏格公式不能直接用于原来的结点三次曲线。

代入微分得
\[
 \frac{\mathrm{d}x}{y}=\frac{2\,\mathrm{d}t}{t^2-1}.
\]
在 \(t=1,-1\) 分别有简单极点，留数为 \(1,-1\)，其余有限点系数均非零。无穷远令 \(s=1/t\)，形式变为 \(-2\,\mathrm{d}s/(1-s^2)\)，全纯且不为0。因此形式没有零点，除子为 \(-[1]-[-1]\)。两个相反留数相加为0，但并不各自为0；结点在平面上的一个位置，归一化后变成了两个需要分别检查的点。原表达在这两点均有极点，因而不是光滑模型上的全纯形式。
\end{solution}
